\documentclass[11pt,a4paper]{article}

\usepackage[utf8]{inputenc}
\usepackage[T1]{fontenc}
\usepackage[french]{babel}
\usepackage{amsmath}
\usepackage{amssymb}
\usepackage{booktabs}
\usepackage{listings}
\usepackage{xcolor}
\usepackage{graphicx}

% Configuration pour les blocs de code VHDL et Verilog
\lstset{
    basicstyle=\ttfamily\footnotesize,
    keywordstyle=\color{blue}\bfseries,
    commentstyle=\color{green!60!black},
    stringstyle=\color{red},
    showstringspaces=false,
    frame=single,
    numbers=left,
    numberstyle=\tiny\color{gray},
    breaklines=true
}

\title{TD1 Correction \\ \large SN360/SN361 : Introduction à la conception des circuits numériques}
\author{Grenoble INP - Esisar - E13}
\date{Année 2024/2025}

\begin{document}

\maketitle

\section*{EXERCICE 1}

\paragraph{1) Rappeler les bornes minimale et maximale pour le codage des nombres signés et non signés sur $n$ bits.}
Pour un nombre signé sur $n$ bits, les valeurs vont de $-2^{(n-1)}$ à $2^{(n-1)}-1$[cite: 6].\\
Pour un nombre non signé sur $n$ bits, les valeurs vont de $0$ à $2^{n}-1$[cite: 7].

\paragraph{2) et 3) Conversion des nombres en base 10 sur 10 bits et hexadécimale :}
Avec 10 bits, on peut coder des nombres signés allant de $-512$ à $511$[cite: 9].\\
Avec 10 bits, on peut coder des nombres non signés allant de $0$ à $1023$[cite: 9].

\begin{table}[h!]
    \centering
    \begin{tabular}{lllc}
        \toprule
        \textbf{Décimal} & \textbf{Signé} & \textbf{Non signé} & \textbf{Hexadécimale} \\
        \midrule
        344 & Oui, 0101011000 & Oui, 0101011000 & 158 \\
        115 & Oui, 0001110011 & Oui, 0001110011 & 73 \\
        -42 & Oui, 1111010110 & Non & FD6 \\
        666 & Non (hors de la plage) & Oui, 1010011010 & 29A \\
        -950 & Non (hors de la plage) & Non & X \\
        -496 & Oui, 1000010000 & Non & E10 \\
        260 & Oui, 0100000100 & Oui, 0100000100 & 104 \\
        1515 & Non (hors de la plage) & Non & X \\
        \bottomrule
    \end{tabular}
\end{table}
% [cite: 10]

\paragraph{4) Pour chaque nombre suivant en base 2 sur 8 bits signés, donner sa valeur numérique correspondante en base 10 :}
On utilise la méthode du complément à 2[cite: 12]:
\begin{itemize}
    \item 0b11000011 $= -61$ [cite: 13]
    \item 0b10100110 $= -90$ [cite: 14]
    \item 0b11110000 $= -16$ [cite: 14]
    \item 0b11111101 $= -3$ [cite: 14]
    \item 0b00111011 $= 59$ [cite: 14]
\end{itemize}

\paragraph{5) En VHDL et en Verilog, comment convertit-on les nombres décimaux 293993 et -293900 en nombres binaires signés et non signés ?}
\textbf{VHDL :}
\begin{lstlisting}[language=VHDL]
library IEEE;
use IEEE.STD_LOGIC_1164.ALL;
use IEEE.NUMERIC_STD.ALL;

entity Convertisseur is
Port (
    resultat_non_signe: out std_logic_vector(18 downto 0);
    resultat_signe: out std_logic_vector(19 downto 0)
);
end Convertisseur;

architecture Behavioral of Convertisseur is
    signal decimal_value: integer := 293993;
    signal binary_value_unsigned: std_logic_vector(18 downto 0);
    signal binary_value_signed: std_logic_vector(19 downto 0);
begin
    binary_value_unsigned <= std_logic_vector(to_unsigned(decimal_value, 19));
    binary_value_signed <= std_logic_vector(to_signed(decimal_value, 20));
    resultat_non_signe <= binary_value_unsigned;
    resultat_signe <= binary_value_signed;
end Behavioral;
\end{lstlisting}
% [cite: 16, 17, 18, 19, 20, 21, 23, 24, 25, 26, 27, 28, 29, 30, 33, 34, 35, 36, 37]

\textbf{Verilog :}
\begin{lstlisting}[language=Verilog]
module Convertisseur (
    output reg [18:0] resultat_non_signe,
    output reg [19:0] resultat_signe
);
    integer decimal_value = 293993;
    always @(*) begin
        resultat_non_signe = decimal_value[18:0];
        resultat_signe = decimal_value[19:0];
    end
endmodule
\end{lstlisting}
% [cite: 38, 39, 40, 41, 42, 43, 44, 45, 46, 48]

\section*{EXERCICE 2}

\paragraph{1) Rappeler la table de vérité de la fonction NAND.}
\begin{table}[h!]
    \centering
    \begin{tabular}{cc|c}
        \toprule
        $a$ & $b$ & $a \text{ NAND } b$ \\
        \midrule
        0 & 0 & 1 \\
        0 & 1 & 1 \\
        1 & 0 & 1 \\
        1 & 1 & 0 \\
        \bottomrule
    \end{tabular}
\end{table}
% [cite: 50, 51, 52, 53, 54, 55, 56, 57, 58, 59, 61, 62, 63, 64, 65]

\paragraph{2) Démontrer par les tables de vérité l'égalité suivante : $b = \overline{a}\cdot b + a\cdot\overline{b}$}
\begin{table}[h!]
    \centering
    \begin{tabular}{cc|cc|cc|c|c}
        \toprule
        $a$ & $b$ & $\overline{a}$ & $\overline{b}$ & $\overline{a}\cdot b$ & $a\cdot\overline{b}$ & $\overline{a}\cdot b + a\cdot\overline{b}$ & $a \text{ XOR } b$ \\
        \midrule
        0 & 0 & 1 & 1 & 0 & 0 & 0 & 0 \\
        0 & 1 & 1 & 0 & 1 & 0 & 1 & 1 \\
        1 & 0 & 0 & 1 & 0 & 1 & 1 & 1 \\
        1 & 1 & 0 & 0 & 0 & 0 & 0 & 0 \\
        \bottomrule
    \end{tabular}
\end{table}
% [cite: 60, 67]

\paragraph{3) En déduire l'expression et un circuit à base de AND, OR et NOT de la fonction XNOR}
En utilisant le tableau de Karnaugh basé sur la table de vérité de l'opération XNOR : $(a\cdot b)+(\overline{a}\cdot\overline{b})$[cite: 69].\\
Remarque : Il est également possible d'appliquer la loi de Morgan à l'équation logique de l'opération XOR : $(a+b)\cdot(\overline{a}+\overline{b})$[cite: 70].

Pour le circuit, on peut réaliser l'équation dérivée du tableau de Karnaugh[cite: 71].
% Espace pour insérer l'image du circuit de la question 3
% \begin{figure}[h!]
%     \centering
%     \includegraphics[width=0.6\textwidth]{chemin/vers/image_circuit_xnor.svg}
% \end{figure}

\paragraph{4) Donner les tables de vérités des fonctions suivantes : $a+a\cdot b$, $a+\overline{a}\cdot b$, multiplexeur 4 vers 1.}

\begin{table}[h!]
    \centering
    \begin{tabular}{cc|c}
        \toprule
        $a$ & $b$ & $a+a\cdot b$ \\
        \midrule
        0 & 0 & 0 \\
        0 & 1 & 0 \\
        1 & 0 & 1 \\
        1 & 1 & 1 \\
        \bottomrule
    \end{tabular}
\end{table}
% [cite: 78]

\begin{table}[h!]
    \centering
    \begin{tabular}{cc|cc|c}
        \toprule
        $a$ & $b$ & $\overline{a}$ & $\overline{a}\cdot b$ & $a+\overline{a}\cdot b$ \\
        \midrule
        0 & 0 & 1 & 0 & 0 \\
        0 & 1 & 1 & 1 & 1 \\
        1 & 0 & 0 & 0 & 1 \\
        1 & 1 & 0 & 0 & 1 \\
        \bottomrule
    \end{tabular}
\end{table}
% [cite: 79]

\begin{table}[h!]
    \centering
    \begin{tabular}{cc|c}
        \toprule
        Sel1 & Sel0 & Sortie \\
        \midrule
        0 & 0 & In0 \\
        0 & 1 & In1 \\
        1 & 0 & In2 \\
        1 & 1 & In3 \\
        \bottomrule
    \end{tabular}
\end{table}
% [cite: 80, 81, 82, 83, 84, 85, 86, 87, 88, 89, 90, 91, 92, 93, 94]

\section*{EXERCICE 3}

\paragraph{1) Quelle fonction logique booléenne retourne 0 lorsque les entrées sont différentes et 1 sinon.}
La fonction est XNOR[cite: 98].

\paragraph{2) Combien de lignes comporte la table de vérité complète ?}
Considérons les valeurs suivantes qui sont sur 4 bits : $A = A_3A_2A_1A_0$ et $B = B_3B_2B_1B_0$. La taille de la table de vérité aura donc $2^{4+4}=256$ lignes[cite: 100].

\paragraph{3) Donner la fonction logique correspondante à la comparaison sur 4 bits.}
Pour que les deux valeurs soient similaires, il faut que chaque bit à l'indice $n$ de A soit égal au bit à l'indice $n$ de B[cite: 102].\\
La fonction logique est : $\text{Comparateur} = (A_3 \text{ XNOR } B_3) \text{ AND } (A_2 \text{ XNOR } B_2) \text{ AND } (A_1 \text{ XNOR } B_1) \text{ AND } (A_0 \text{ XNOR } B_0)$[cite: 104].

\paragraph{4) Faire le schéma de câblage}
% Espace pour insérer l'image du circuit comparateur 4 bits
% \begin{figure}[h!]
%     \centering
%     \includegraphics[width=0.8\textwidth]{chemin/vers/image_comparateur.svg}
% \end{figure}
[cite: 105]

\paragraph{5) Ecrire cette fonction en VHDL et Verilog comportemental.}
\textbf{VHDL :}
\begin{lstlisting}[language=VHDL]
library IEEE;
use IEEE.STD_LOGIC_1164.ALL;

entity Comparateur_4bits is
Port (
    A: in STD_LOGIC_VECTOR(3 downto 0);
    B: in STD_LOGIC_VECTOR(3 downto 0);
    Equal: out STD_LOGIC
);
end Comparateur_4bits;

architecture Behavioral of Comparateur_4bits is
begin
    Equal <= (A(3) xnor B(3)) and (A(2) xnor B(2)) and (A(1) xnor B(1)) and (A(0) xnor B(0));
end Behavioral;
\end{lstlisting}
% [cite: 111, 112, 113, 114, 115, 117, 118, 119, 120, 121, 122, 123]

\textbf{Verilog :}
\begin{lstlisting}[language=Verilog]
module Comparateur_4bits (
    input [3:0] A,
    input [3:0] B,
    output Equal
);
    assign Equal = (A[3] ~^ B[3]) & (A[2] ~^ B[2]) & (A[1] ~^ B[1]) & (A[0] ~^ B[0]);
endmodule
\end{lstlisting}
% [cite: 125, 126, 128, 129, 130, 131]

\section*{EXERCICE 4}

Eq1 : $a\cdot b + \overline{a}\cdot\overline{b}$ [cite: 133]\\
Eq2 : $(a+b)\cdot c$ [cite: 134]\\
Eq3 : $\overline{a\cdot b}+c$ [cite: 135]

\paragraph{1) Représentez avec des portes logiques les trois équations.}
[cite: 137]

\paragraph{2) Codez en VHDL et Verilog les équations en comportemental.}
On le fait pour l'équation 1[cite: 140]:

\textbf{VHDL :}
\begin{lstlisting}[language=VHDL]
library IEEE;
use IEEE.STD_LOGIC_1164.ALL;

entity Eq1 is
Port (
    a: in STD_LOGIC;
    b: in STD_LOGIC;
    result: out STD_LOGIC
);
end Eq1;

architecture Behavioral of Eq1 is
begin
    result <= (a AND b) OR (NOT a AND NOT b);
end Behavioral;
\end{lstlisting}
% [cite: 141, 142, 143, 147, 149, 150, 151, 152, 153, 154, 155, 156]

\textbf{Verilog :}
\begin{lstlisting}[language=Verilog]
module Eq1 (
    input wire a,
    input wire b,
    output wire result
);
    assign result = (a & b) | (~a & ~b);
endmodule
\end{lstlisting}
% [cite: 157, 159, 160, 161, 162, 163, 164]

\section*{EXERCICE 5}

\paragraph{1) Compléter la table de vérité de l'afficheur 7 segments.}

\begin{table}[h!]
    \centering
    \begin{tabular}{ccc}
        \toprule
        \textbf{Decimal Digit} & \textbf{Input lines (ABCD)} & \textbf{Output lines (abcdefg)} \\
        \midrule
        0 & 0000 & 1111110 \\
        1 & 0001 & 0110000 \\
        2 & 0010 & 1101101 \\
        3 & 0011 & 1111001 \\
        4 & 0100 & 0110011 \\
        5 & 0101 & 1011011 \\
        6 & 0110 & 1011111 \\
        7 & 0111 & 1110000 \\
        8 & 1000 & 1111111 \\
        9 & 1001 & 1111011 \\
        \bottomrule
    \end{tabular}
\end{table}
% [cite: 170, 171, 172, 173, 174, 175, 176, 177, 178, 179, 181, 182, 184, 186, 188, 189, 191, 192, 194, 199, 200, 201, 203, 204]

\paragraph{2) En utilisant les tableaux de Karnaugh, simplifier les fonctions des segments $a$, $b$ et $g$.}
[cite: 206]
\begin{align*}
    a &= A + C + B\cdot D + \overline{B}\cdot\overline{D} \\
    b &= \overline{B} + \overline{C}\cdot\overline{D} + C\cdot D \\
    g &= A + B\cdot\overline{C} + \overline{B}\cdot C + C\cdot\overline{D}
\end{align*}
% [cite: 275, 276, 277]

\end{document}